% AI-produced modern German version of the AI-transcribed source component: printed 162--171.
\begingroup
\clearpage
\noindent162\hfill\textsc{Nachlass.}\hfill\mbox{}\par\medskip
\begin{center}
\begingroup\fontsize{7.6}{9}\selectfont\setlength{\tabcolsep}{2pt}\renewcommand{\arraystretch}{1.07}
\begin{tabular}{|r|r|r|r|r||r|r|r|r|r|}
\hline
\rule{0pt}{18pt}\shortstack{$x$\\(Meter)}&$\varphi$&\shortstack{$\log\alpha$\\$1{,}....-10$}&\shortstack{$\log\beta$\\$7{,}....-30$}&\shortstack{$\log\gamma$\\$2{,}....-20$}&\shortstack{$x$\\(Meter)}&$\varphi$&\shortstack{$\log\alpha$\\$1{,}....-10$}&\shortstack{$\log\beta$\\$7{,}....-30$}&\shortstack{$\log\gamma$\\$2{,}....-20$}\\\hline
$-200\,000$&$53^\circ19'40''{,}521$&$1{,}4033800$&$7{,}31422$&$2{,}40145$&$-500\,000$&$56^\circ1'24''{,}643$&$1{,}4031246$&$7{,}31397$&$2{,}38676$\\
$210\,000$&$25\ 4{,}061$&$713$&$22$&$103$&$510\,000$&$6\ 48{,}041$&$163$&$96$&$620$\\
$220\,000$&$30\ 27{,}597$&$627$&$21$&$060$&$520\,000$&$12\ 11{,}434$&$1{,}4031079$&$95$&$563$\\
$230\,000$&$35\ 51{,}127$&$541$&$20$&$2{,}40017$&$530\,000$&$17\ 34{,}822$&$1{,}4030996$&$94$&$506$\\
$240\,000$&$41\ 14{,}653$&$455$&$19$&$2{,}39974$&$540\,000$&$22\ 58{,}205$&$912$&$93$&$448$\\
\hline
$-250\,000$&$46\ 38{,}173$&$369$&$18$&$930$&$-550\,000$&$28\ 21{,}584$&$829$&$92$&$389$\\
$260\,000$&$52\ 1{,}689$&$282$&$17$&$885$&$560\,000$&$33\ 44{,}958$&$746$&$92$&$330$\\
$270\,000$&$53\ 57\ 25{,}200$&$197$&$16$&$840$&$570\,000$&$39\ 8{,}329$&$663$&$91$&$271$\\
$280\,000$&$54\ 2\ 48{,}707$&$111$&$15$&$795$&$580\,000$&$44\ 31{,}693$&$580$&$90$&$211$\\
$290\,000$&$8\ 12{,}208$&$1{,}4033025$&$15$&$749$&$590\,000$&$49\ 55{,}053$&$498$&$\raisebox{.6ex}{\tiny\textquoteleft}89$&$150$\\
\hline
$-300\,000$&$13\ 35{,}705$&$1{,}4032939$&$14$&$703$&$-600\,000$&$56\ 55\ 18{,}409$&$415$&$88$&$089$\\
$310\,000$&$18\ 59{,}197$&$854$&$13$&$656$&$610\,000$&$57\ 0\ 41{,}760$&$333$&$87$&$2{,}38028$\\
$320\,000$&$24\ 22{,}684$&$768$&$12$&$609$&$620\,000$&$6\ 5{,}106$&$250$&$87$&$2{,}37966$\\
$330\,000$&$29\ 46{,}167$&$683$&$11$&$561$&$630\,000$&$11\ 28{,}448$&$168$&$86$&$904$\\
$340\,000$&$35\ 9{,}644$&$598$&$10$&$513$&$640\,000$&$16\ 51{,}785$&$086$&$85$&$841$\\
\hline
$-350\,000$&$40\ 33{,}117$&$512$&$09$&$464$&$-650\,000$&$22\ 15{,}118$&$1{,}4030004$&$84$&$777$\\
$360\,000$&$45\ 56{,}585$&$427$&$09$&$415$&$660\,000$&$27\ 38{,}446$&$1{,}4029922$&$83$&$713$\\
$370\,000$&$51\ 20{,}049$&$342$&$08$&$366$&$670\,000$&$33\ 1{,}770$&$840$&$83$&$648$\\
$380\,000$&$54\ 56\ 43{,}508$&$257$&$07$&$316$&$680\,000$&$38\ 25{,}089$&$759$&$82$&$583$\\
$390\,000$&$55\ 2\ 6{,}961$&$173$&$06$&$265$&$690\,000$&$43\ 48{,}403$&$677$&$81$&$518$\\
\hline
$-400\,000$&$7\ 30{,}411$&$088$&$05$&$214$&$-700\,000$&$49\ 11{,}713$&$596$&$80$&$452$\\
$410\,000$&$12\ 53{,}855$&$1{,}4032003$&$04$&$162$&$710\,000$&$54\ 35{,}018$&$514$&$79$&$385$\\
$420\,000$&$18\ 17{,}295$&$1{,}4031919$&$03$&$110$&$720\,000$&$57\ 59\ 58{,}319$&$433$&$78$&$318$\\
$430\,000$&$23\ 40{,}730$&$834$&$03$&$058$&$730\,000$&$58\ 5\ 21{,}615$&$352$&$78$&$250$\\
$440\,000$&$29\ 4{,}160$&$750$&$02$&$2{,}39005$&$740\,000$&$10\ 44{,}907$&$271$&$77$&$182$\\
\hline
$-450\,000$&$34\ 27{,}586$&$666$&$01$&$2{,}38951$&$-750\,000$&$16\ 8{,}194$&$190$&$76$&$113$\\
$460\,000$&$39\ 51{,}007$&$582$&$7{,}31400$&$897$&$760\,000$&$21\ 31{,}476$&$110$&$75$&$2{,}37044$\\
$470\,000$&$45\ 14{,}423$&$498$&$7{,}31399$&$843$&$770\,000$&$26\ 54{,}755$&$1{,}4029029$&$74$&$2{,}36974$\\
$480\,000$&$50\ 37{,}834$&$414$&$98$&$788$&$780\,000$&$32\ 18{,}028$&$1{,}4028949$&$73$&$904$\\
$490\,000$&$55\ 56\ 1{,}241$&$330$&$98$&$732$&$790\,000$&$37\ 41{,}297$&$869$&$73$&$833$\\
\hline
$-500\,000$&$56\ 1\ 24{,}643$&$1{,}4031246$&$7{,}31397$&$2{,}38676$&$-800\,000$&$58\ 43\ 4{,}562$&$1{,}4028788$&$7{,}31372$&$2{,}36762$\\
\end{tabular}\par\endgroup
\end{center}
[Dieser Tabelle liegt der Abplattungswert $\frac1{302{,}68}$ zugrunde. Die Werte von $x$ sind um $-26{,}264\,\mathrm m$ zu korrigieren, wenn für Göttingen $\varphi=51^\circ31'47''{,}85$ statt $51^\circ31'48''{,}70$ angenommen wird.]
\bigskip
\begin{center}
[Der Unterschied zwischen der Projektion der geodätischen Linie\\und der ihre Endpunkte verbindenden Geraden bei der konformen Abbildung\\einer gekrümmten Fläche in der Ebene.]\par\medskip
[16.]
\end{center}
Ein sehr fruchtbares Prinzip für die allgemeine Theorie der gekrümmten Flächen ist das folgende.

Die einzelnen Punkte der gekrümmten Fläche werden auf beliebige Weise durch zwei veränderliche Größen $p$, $q$ [dargestellt]. Alle Punkte,
\clearpage
\noindent\mbox{}\hfill\textsc{Konforme Abbildung des Sphäroids in der Ebene.}\hfill163\par\medskip
\noindent für die $q$ konstant ist, bilden also auf der Fläche eine Linie. Ihr Element sei $\mathrm du=t\,\mathrm dp$, wobei $t$ eine Funktion von $p$ und $q$ sein wird.

Zu dem beliebigen Punkt der gekrümmten Fläche, der durch $p$, $q$ bestimmt wird, werde nun eine kürzeste Linie gezogen, entweder von einem festen Punkt aus oder senkrecht zu einer gegebenen Linie. Ihre Länge $r$ wird also selbst eine Funktion von $p$ und $q$ sein. Ist nun $\varphi$ der Winkel, den das Element von $r$ mit dem Element $\mathrm du$ bildet, so gilt
\[
\frac{\partial r}{\partial p}=t\cos\varphi.
\]
Eine kürzeste Linie auf der gekrümmten Fläche werde nun bei der elementähnlichen Abbildung in der Ebene durch die Kurve dargestellt, deren Punkte durch $R$, den Abstand vom Anfangspunkt, [und] $\theta$, die Richtung, bestimmt werden. Das [reziproke] Vergrößerungsverhältnis sei $n^*$, d.~h., einander entsprechende Elemente in der Ebene und auf der Fläche verhalten sich wie $1$ zu $n^*$.

Es sei ferner
\[
\begin{aligned}
n^*&=N+aR+bRR+cR^3+\ldots\\
\rlap{\hspace{-1.4em}\raisebox{.7ex}{.}}\log n^*&=\log N+\alpha R+\beta RR+\gamma R^3+\ldots,
\end{aligned}
\]
wobei $a$, $b$, $c$, $\ldots$ und $\alpha$, $\beta$, $\gamma$, $\ldots$ Funktionen von $\theta$ bezeichnen.

Man hat dann [$\mathrm dn^*=n^*\mathrm d\log n^*$ oder]
\[
(a+2bR+3cRR\ldots)=(N+aR+bRR\ldots)(\alpha+2\beta R+3\gamma RR\ldots)
\]
[und hieraus]
\[
\begin{aligned}
a&=N\alpha\\
2b&=2N\beta+a\alpha\\
3c&=3N\gamma+2a\beta+b\alpha
\end{aligned}
\]
\begin{center}etc.\end{center}
[oder]
\[
\begin{aligned}
a&=N\alpha\\
b&=N\beta+\tfrac12N\alpha\alpha\\
c&=N\gamma+N\alpha\beta+\tfrac16N\alpha^3
\end{aligned}
\]
\begin{center}etc.\end{center}
\vfill\hfill$21^*$\hspace*{2cm}\par
\clearpage
\noindent164\hfill\textsc{Nachlass.}\hfill\mbox{}\par\medskip
\begingroup\small
Man hat nun nach obigem
\[
\begin{gathered}
\frac{\partial r}{\partial R}=n^*\cos\varphi\\
\frac{\partial r}{R\partial\theta}=-n^*\sin\varphi,\\
\left(\frac{\partial r}{\partial R}\right)^2+\left(\frac{\partial r}{R\partial\theta}\right)^2={n^*}^2;
\end{gathered}
\]
[ferner ist
\[
m\frac{\partial(\theta+\psi)}{\partial R}=n^*\sin\varphi,\qquad m\frac{\partial(\theta+\psi)}{R\partial\theta}=n^*\cos\varphi,
\]
wobei $m\,\mathrm d(\theta+\psi)$ das Element der Kurve mit konstantem $r$ bezeichnet. Hieraus folgt
\[
\frac{\partial(\theta+\psi)}{\partial R}\cos\varphi-\frac{\partial(\theta+\psi)}{R\partial\theta}\sin\varphi=0
\]
oder]
\[
\frac{\partial(\theta+\psi)}{\partial R}.\frac{\partial r}{\partial R}+\frac{\partial(\theta+\psi)}{R\partial\theta}.\frac{\partial r}{R\partial\theta}=0.
\]
[$\psi$ und $\varphi$ bezeichnen die Winkel, die das Bild der geodätischen Linie in der Ebene mit $R$ am Anfangs- bzw. am Endpunkt bildet.]

Man findet hieraus, wenn man
\[
\frac{\partial a}{\partial\theta}=a',\qquad\frac{\partial b}{\partial\theta}=b',\ \text{usw.;}\qquad
\frac{\partial a'}{\partial\theta}=a'',\qquad\frac{\partial b'}{\partial\theta}=b'',\ \text{usw.}
\]
setzt und bemerkt, dass $a''=-a$ [ist, weil $a$ die Form $C_1\cos\theta+C_2\sin\theta$ hat]:
\[
\begin{aligned}
r&=NR+\tfrac12aRR+\left(\tfrac13b-\frac{a'a'}{24N}\right)R^3+\left(\tfrac14c-\frac{a'b'}{24N}+\frac{aa'a'}{48NN}\right)R^4\ldots;\\
n^*\cos\varphi&=N+aR+\left(b-\frac{a'a'}{8N}\right)RR+\left(c-\frac{a'b'}{6N}+\frac{aa'a'}{12NN}\right)R^3\ldots\\
-n^*\sin\varphi&=\tfrac12a'R+\left(\tfrac13b'+\frac{aa'}{12N}\right)RR
+\left(\tfrac14c'+\frac{ab'}{24N}-\frac{a'b''}{24N}+\frac{{a'}^3}{48NN}-\frac{aaa'}{24NN}\right)R^3\ldots.
\end{aligned}
\]
Also
\[
\begin{aligned}
-\sin\varphi&=\frac{a'}{2N}R+\left(\frac{b'}{3N}-\frac{5aa'}{12NN}\right)RR\ldots\\
-\varphi&=\frac{a'}{2N}R+\left(\frac{b'}{3N}-\frac{5aa'}{12NN}\right)RR\ldots\\
&=\tfrac12\alpha'R+\left(\tfrac13\beta'-\tfrac1{12}\alpha\alpha'\right)RR\ldots.
\end{aligned}
\]
Es findet sich ferner aus obiger partieller Differentialgleichung
\[
\begin{aligned}
-\psi&=\frac{a'}{2N}R+\left(\frac{b'}{6N}-\frac{aa'}{12NN}\right)RR\ldots\\
&=\tfrac12\alpha'R+\left(\tfrac16\beta'+\tfrac1{12}\alpha\alpha'\right)RR\ldots.
\end{aligned}
\]
Wenn die Lage der Punkte in der Ebene durch $x$, $y$ ausgedrückt wird, wobei $x=R\cos\theta$, $y=R\sin\theta$, sei
\par\endgroup
\clearpage
\noindent\mbox{}\hfill\textsc{Konforme Abbildung des Sphäroids in der Ebene.}\hfill165\par\medskip
\begingroup\small
\[
\log n^*=\log N-Ax-By-\tfrac12Cxx-Dxy-\tfrac12Eyy\ldots,
\]
daraus folgt
\[
\begin{array}{rl@{\qquad}rl}
\alpha R&=-Ax-By&\alpha'R&=Ay-Bx\\
\beta RR&=-\tfrac12Cxx-Dxy-\tfrac12Eyy&\beta'RR&=(C-E)xy-D(xx-yy)\\
&\text{etc.}&&\text{etc.}
\end{array}
\]
\[
\begin{aligned}
\varphi&=+\tfrac12Bx+(\tfrac13D+\tfrac1{12}AB)(xx-yy)\\
&\qquad-\tfrac12Ay+(-\tfrac13C+\tfrac13E-\tfrac1{12}AA+\tfrac1{12}BB)xy\ldots\\
\psi&=+\tfrac12Bx+(\tfrac16D-\tfrac1{12}AB)(xx-yy)\\
&\qquad-\tfrac12Ay+(-\tfrac16C+\tfrac16E+\tfrac1{12}AA-\tfrac1{12}BB)xy\ldots;\\[5pt]
r&=N\left(R+\tfrac12\alpha RR+(\tfrac13\beta+\tfrac16\alpha\alpha-\tfrac1{24}\alpha'\alpha')R^3\ldots\right)\\
\log r&=\log NR+\tfrac12\alpha R+\left(\tfrac13\beta+\tfrac1{24}(\alpha\alpha-\alpha'\alpha')\right)RR\ldots\\
&=\log N^0R+\left(\tfrac1{12}\beta+\tfrac1{24}(\alpha\alpha-\alpha'\alpha')\right)RR\ldots\\
&=\log N^0R-\tfrac1{24}(C+BB-AA)xx\\
&\qquad\qquad-\tfrac1{12}(D-2AB)xy\\
&\qquad\qquad-\tfrac1{24}(E+AA-BB)yy
\end{aligned}
\]
\begin{center}etc.,\end{center}
wobei $N^0$ der Wert von $n^*$ für den Mittelpunkt der Geraden ist [also $\log N^0=\log N+\tfrac12\alpha R+\tfrac14\beta RR+\ldots$ gilt].

Es sei nun
\[
\begin{aligned}
\mathopen{[}\log n=]\log\frac1{n^*}&=(e+e'x+e''xx+\ldots)(Y+y)^2\\
&\quad+(f+f'x+f''xx+\ldots)(Y+y)^4\\
&\quad+(g+g'x+g''xx+\ldots)(Y+y)^6
\end{aligned}
\]
\begin{center}etc.,\end{center}
wobei die zweite Koordinate [am Anfangspunkt] $=Y$ [ist]. Dann gilt
\[
\begin{aligned}
-\log N&=eYY+fY^4+gY^6\ldots\\
A&=e'YY+f'Y^4+g'Y^6\ldots\\
B&=2eY+4fY^3+6gY^5\ldots\\
\tfrac12C&=e''YY+f''Y^4+g''Y^6\ldots\\
D&=2e'Y+4f'Y^3+6g'Y^5\ldots\\
\tfrac12E&=e+6fYY+15gY^4\ldots
\end{aligned}
\]
\begin{center}etc.\end{center}
\par\endgroup
\clearpage
\noindent166\hfill\textsc{Nachlass.}\hfill\mbox{}\par\medskip
Bei der Abbildung der Kugelfläche mit dem Radius $h$ ist $n^*$ von $x$ unabhängig, und
\[
\begin{aligned}
\frac1{n^*}&=1+\tfrac12\frac{(Y+y)^2}{hh}+\tfrac1{24}\frac{(Y+y)^4}{h^4}\ldots\\
\log\frac1{n^*}&=\tfrac12\frac{(Y+y)^2}{hh}-\tfrac1{12}\frac{(Y+y)^4}{h^4}\ldots.
\end{aligned}
\]
Also [ist]
\[
A=C=D=0
\]
[und]
\[
B=\frac{Y}{hh}-\tfrac13\frac{Y^3}{h^4}\ldots.
\]
Das Glied $\tfrac1{12}BBxy$ [in $\varphi$ und $\psi$] kann also nicht größer als $\tfrac1{24}\frac{YY}{h^4}RR$ werden, also für $R=\lambda.100000\,\mathrm m$, $Y=\mu.400000\,\mathrm m$ [nicht größer als] $\lambda\lambda.\mu\mu.0''{,}008$.

Für $\varphi$ und $\psi$ werden die Glieder [auch beim Sphäroid im Wesentlichen] durch die Formel $\tfrac12Bx-\tfrac12Ay$ erfasst, wenn man in den beiden Fällen die Werte von $B$ und $A$ an den Punkten der Geraden nimmt, die um $\frac23$ bzw. $\frac13$ ihrer Länge vom Anfangspunkt entfernt sind. Diese Werte sind oben [Art.~15] mit hinreichender Genauigkeit angegeben.

[Sind $A_{\frac23}$, $B_{\frac23}$ und $A_{\frac13}$, $B_{\frac13}$ die gerade bezeichneten Werte von $A$ und $B$, so gilt
\[
\begin{aligned}
\varphi&=\tfrac12B_{\frac23}.x-\tfrac12A_{\frac23}.y+\tfrac13(Dx-Cy)x+\tfrac1{12}(Ax+By)(Bx-Ay)\ldots\\
\psi&=\tfrac12B_{\frac13}.x-\tfrac12A_{\frac13}.y+\tfrac16(Dx-Cy)x+\tfrac1{12}(Ax+By)(-Bx+Ay)\ldots.]
\end{aligned}
\]
\begin{center}[17.]\end{center}
Die Verallgemeinerung der obigen Grundsätze führt zu folgender Behandlung.

Es seien $p$, $q$ zwei veränderliche Größen, die die verschiedenen Punkte der gekrümmten Fläche bestimmen. Die Länge der kürzesten Linie heiße $r$, ihre Richtung am Anfang werde durch $\psi$ bezeichnet. Schließlich sei $\int\rho\,\mathrm d\psi$ die Länge der Linie mit konstantem $r$ und $\mathrm ds$ [das] Element einer Linie auf der gekrümmten Fläche, das allgemein durch $\sqrt{(A\,\mathrm dp^2+2B\,\mathrm dp\mathrm dq+C\,\mathrm dq^2)}$ ausgedrückt werde.

Man setze
\[
\begin{aligned}
\mathrm dr&=g\,\mathrm dp+h\,\mathrm dq\\
\rho\,\mathrm d\psi&=G\,\mathrm dp+H\,\mathrm dq.
\end{aligned}
\]
\clearpage
\noindent\mbox{}\hfill\textsc{Konforme Abbildung des Sphäroids in der Ebene.}\hfill167\par\medskip
Hieraus [folgt]:
\[
\begin{aligned}
gg+GG&=A\\
gh+GH&=B\\
hh+HH&=C,
\end{aligned}
\]
also [wird]
\[
(A-gg)(C-hh)=(B-gh)^2
\]
[oder]
\[
\begin{array}{r@{\quad}l}
\mathrm I.&\begin{aligned}[t]
AC-BB&=Ahh-2Bgh+Cgg\\
&\left[=A\left(\frac{\partial r}{\partial q}\right)^2-2B\frac{\partial r}{\partial p}\frac{\partial r}{\partial q}+C\left(\frac{\partial r}{\partial p}\right)^2.\right.
\end{aligned}
\end{array}
\]
Ferner ist]
\[
\begin{aligned}
Ah-Bg&=G(Gh-Hg)\\
Bh-Cg&=H(Gh-Hg),
\end{aligned}
\]
also
\[
\frac G\rho(Cg-Bh)+\frac H\rho(Ah-Bg)=0,
\]
d.~h.
\[
\begin{array}{r@{\qquad}l}
\mathrm{II}.&\displaystyle\frac{\partial\psi}{\partial p}(Cg-Bh)+\frac{\partial\psi}{\partial q}(Ah-Bg)=0.
\end{array}
\]
Bilden die Linien mit konstantem $p$ bzw. $q$ mit den Linien mit konstantem $\psi$ die Winkel $\mathrm M$ bzw. $\mathrm N$, so gilt
\[
\begin{array}{rl@{\qquad}rl}
g&=\cos\mathrm N\sqrt A&[G&=\sin\mathrm N\sqrt A\\
h&=\cos\mathrm M\sqrt C&H&=\sin\mathrm M\sqrt C;
\end{array}
\]
\rlap{\hspace{1.2em}\raisebox{-1.1ex}{.}}mithin]
\[
B=\cos(\mathrm M-\mathrm N)\sqrt{AC}.
\]
Wenn z.~B.
\[
\begin{array}{rl@{\qquad}rl}
p&=R&x&=R\cos\theta\\
q&=\theta&y&=R\sin\theta\\
&&z&=1+tRR+uR^3\ldots
\end{array}
\]
ist, wobei $t$, $u$ Funktionen von $\theta$ sind, [so gilt]
\[
\mathrm ds=\sqrt{\{\mathrm dR^2+RR\,\mathrm d\theta^2+\mathrm dz^2\}},
\]
und
\[
A=1+\left(\frac{\partial z}{\partial R}\right)^2,\qquad
B=\frac{\partial z}{\partial R}.\frac{\partial z}{\partial\theta},\qquad
C=RR+\left(\frac{\partial z}{\partial\theta}\right)^2.
\]
[Nach Gleichung I gilt dann, wenn $t'$, $u'$, $\ldots$ die Ableitungen von $t$, $u$, $\ldots$ nach $\theta$ sind:
\clearpage
\noindent168\hfill\textsc{Nachlass.}\hfill\mbox{}\par\medskip
\begingroup\small
\[
\begin{aligned}
&(1+t't'RR+2t'u'R^3\ldots)\left(\frac{\partial r}{\partial R}\right)^2
+(1+4ttRR+12tuR^3\ldots)\left(\frac{\partial r}{R\partial\theta}\right)^2\\
&\qquad-2(2tt'RR+(2tu'+3ut')R^3\ldots)\frac{\partial r}{\partial R}.\frac{\partial r}{R\partial\theta}\\
&\quad=1+(4tt+t't')RR+(12tu+2t'u')R^3\ldots,
\end{aligned}
\]
woraus sich ergibt:]
\[
r=R\qquad\ast\qquad+\tfrac23ttR^3+\tfrac32tuR^4\ldots.
\]
[Gleichung II ergibt damit:
\[
\begin{aligned}
&(1+(2tt+t't')RR+(6tu+2t'u')R^3\ldots)\frac{\partial\psi}{\partial R}\\
&+(-\tfrac23tt'RR-(\tfrac12tu'+\tfrac32ut')R^3\ldots)\frac{\partial\psi}{R\partial\theta}=0,
\end{aligned}
\]
folglich wird:
\[
\psi=\theta+\tfrac13tt'RR+(\tfrac12ut'+\tfrac16tu')R^3\ldots.]
\]
\bigskip
\begin{center}
Zur Transformation der Koordinaten.\par\medskip
[18.]
\end{center}
Die Grundgleichung ist
\[
\mathrm f(x+iy+\omega)-\tfrac12\mathrm f(x+iy)+\tfrac12\mathrm f(x-iy)=\mathrm f(x+\Omega).
\]
Wir schreiben
\[
\begin{aligned}
\mathrm f'(x)&=a\\
\mathrm f''(x)&=b\\
\mathrm f'''(x)&=c\\
\mathrm f^{\mathrm{IV}}(x)&=d
\end{aligned}
\]
\begin{center}etc.\end{center}
[Dann wird]
\[
\begin{aligned}
&\mathrm f(x)+a\omega+\tfrac12b(\omega\omega+2iy\omega-yy)+\tfrac16c(\omega^3+3iy\omega\omega-3yy\omega)\\
&\qquad+\tfrac1{24}d(\omega^4+4iy\omega^3-6yy\omega\omega-4iy^3\omega+y^4)\ldots\\
&\quad=\mathrm f(x)+a\Omega+\tfrac12b\Omega\Omega+\tfrac16c\Omega^3+\tfrac1{24}d\Omega^4\ldots
\end{aligned}
\]
[und hieraus]
\par\endgroup
\clearpage
\noindent\mbox{}\hfill\textsc{Konforme Abbildung des Sphäroids in der Ebene.}\hfill169\par\medskip
\begingroup\small
\[
\begin{aligned}
x+\Omega&=x-\frac b{2a}yy+\left(\frac d{24a}\mathbin{\overset{\scriptscriptstyle .}{-}}\frac{b^3}{8a^3}\right)y^4\ldots\\
&\quad+\omega\left\{1+\frac ba.iy-\left(\frac c{2a}-\frac{bb}{2aa}\right)yy-\left(\frac d{6a}-\frac{b^3}{2a^3}\right)iy^3\ldots\right\}\\
&\quad+\omega\omega\left\{\left(\frac c{2a}-\frac{bb}{aa}\right)iy-\left(\frac d{4a}-\frac{3bc}{4aa}\right)yy\ldots\right\}\\
&\quad+\omega^3\left\{\left(\frac d{6a}-\frac{bc}{aa}+\frac{b^3}{a^3}\right)iy\ldots\right\}
\end{aligned}
\]
\begin{center}etc.\end{center}
\bigskip
\begin{center}[19.]\end{center}
Die Entwicklung beruht auf folgender Aufgabe.

Wenn
\[
u-Auu-Bu^3-Cu^4-Du^5-\ldots=p+Aq+Br+Cs+Dt+\ldots,
\]
[wo $p$ von derselben Ordnung wie $u$, $q$ von der Ordnung $uu$, usw. ist,] so wird:
\[
\begin{aligned}
u&=p+A(pp+q)+AA(2p^3+2pq)+A^3(5p^4+6ppq+qq)\\
&\qquad\qquad+B(\ p^3+r)\qquad+AB(5p^4+3ppq+2pr)\\
&\qquad\qquad\qquad\qquad\qquad+C(\ p^4+s)\\
&\quad+\phantom{AA}A^4(14p^5+20p^3q+6pqq)\\
&\quad+AAB(21p^5+20p^3q+6ppr+3pqq+2qr)\\
&\quad+\phantom{A}AC(\ 6p^5+\ 4p^3q+2ps)\\
&\quad+\phantom{A}BB(\ 3p^5+3ppr)\\
&\quad+\phantom{AA}D(\ p^5+t)
\end{aligned}
\]
\begin{center}etc.\end{center}
Allgemein
\[
u=\sum A^\alpha B^\beta C^\gamma\ldots p^\pi q^\chi r^\rho s^\sigma\ldots.
\frac{\Pi(\alpha+\beta+\gamma+\ldots+\pi+\chi+\rho+\sigma+\ldots-1)}{\Pi(\alpha-\chi).\Pi(\beta-\rho).\Pi(\gamma-\sigma)\ldots\Pi\pi.\Pi\chi.\Pi\rho.\Pi\sigma\ldots},
\]
wo
\[
\alpha,\ \beta,\ \gamma,\ \ldots,\ \pi,\ \chi,\ \rho,\ \sigma,\ \ldots,\ \alpha-\chi,\ \beta-\rho,\ \gamma-\sigma,\ \ldots
\]
nichtnegative ganze Zahlen sein müssen und
\[
1+\alpha+2\beta+3\gamma+\ldots=\pi+2\chi+3\rho+4\sigma+\ldots
\]
[ist].
\par\endgroup
\vfill\noindent{\footnotesize IX.}\hfill22\hspace*{2cm}\par
\clearpage
\noindent170\hfill\textsc{Nachlass.}\hfill\mbox{}\par\medskip
Man hat nun
\[
u=\Omega
\]
[und]
\[
\begin{aligned}
\rlap{\hspace{-1.3em}.}p&=-z+(\omega+z)=\omega\\
q&=\phantom{z^3}-(\omega+z)^2\\
r&=z^3-(\omega+z)^3\\
s&=\phantom{z^3}-(\omega+z)^4\\
t&=z^5-(\omega+z)^5
\end{aligned}
\]
\begin{center}etc.\end{center}
\noindent\llap{\raisebox{-.4ex}{.}\,}[zu setzen]. Also
\begingroup\small
\[
\begin{aligned}
\Omega&=\omega-A(2\omega z+zz)-AA(4\omega\omega z+2\omega zz)\\
&\qquad\qquad\qquad\qquad-B(3\omega\omega z+3\omega zz)\\[5pt]
&\quad-A^3(\ 8\omega^3z\qquad\ast\qquad-4\omega z^3-z^4)\\
&\quad-AB(12\omega^3z+9\omega\omega zz\qquad\ast\qquad\ast\ )\\
&\quad-C(\ 4\omega^3z+6\omega\omega zz+4\omega z^3+z^4)\\[5pt]
&\qquad\qquad-A^4(16\omega^4z-16\omega^3zz-24\omega\omega z^3-6\omega z^4)\\
&\qquad\qquad-AAB(36\omega^4z\qquad\ast\qquad-30\omega\omega z^3-9\omega z^4)\\
&\qquad\qquad-AC(16\omega^4z+16\omega^3zz+8\omega\omega z^3+2\omega z^4)\\
&\qquad\qquad-BB(\ 9\omega^4z+9\omega^3zz\qquad\ast\qquad\ast\ )\\
&\qquad\qquad-D(\ 5\omega^4z+10\omega^3zz+10\omega\omega z^3+5\omega z^4)
\end{aligned}
\]
\begin{center}etc.\end{center}
\par\endgroup
Für die Übertragung in die vorherige Schreibweise [Art.~18] braucht nur noch gesetzt zu werden:
\[
\begin{aligned}
z&=iy\\[5pt]
A&=-\frac b{2a}\\[5pt]
B&=-\frac c{6a}\\[5pt]
C&=-\frac d{24a}\\[5pt]
D&=-\frac e{120a}
\end{aligned}
\]
\begin{center}etc.\end{center}
\clearpage
\noindent\mbox{}\hfill\textsc{Konforme Abbildung des Sphäroids in der Ebene.}\hfill171\par\medskip
\begin{center}[20.]\end{center}
Auf der Kugel [mit dem Radius $1$ gilt nach S.~144, wenn $x$ von einem beliebigen Anfangspunkt aus in südlicher Richtung positiv gezählt wird,
\[
\frac{\mathrm d\varphi}{\mathrm dx}=-1,\qquad\mathrm f'(x)=-\frac1{\cos\varphi};
\]
setzt man
\[
\sin\varphi=s,\qquad\cos\varphi=c,]\qquad\operatorname{tang}\varphi=t,
\]
[so wird mithin
\begingroup\small
\[
\begin{array}{ll@{\quad}l}
\displaystyle\rlap{\hspace{-.3em}\raisebox{1.2ex}{.}}\mathrm f''(x)=+\frac{s}{cc}&\text{und daher}&A=+\tfrac12t\\[6pt]
\displaystyle\mathrm f'''(x)=-\frac{1+ss}{c^3}&&B=-\tfrac16(1+2tt)\\[6pt]
\displaystyle\mathrm f^{\mathrm{IV}}(x)=+\frac{5s+s^3}{c^4}&&C=+\tfrac1{24}(5t+6t^3)\\[6pt]
\displaystyle\mathrm f^{\mathrm V}(x)=-\frac{5+18ss+s^4}{c^5}&&D=-\tfrac1{120}(5+28tt+24t^4)\\[6pt]
\displaystyle\mathrm f^{\mathrm{VI}}(x)=+\frac{61s+58s^3+s^5}{c^6}&&E=+\tfrac1{720}(61t+180t^3+120t^5)\\[6pt]
\displaystyle\mathrm f^{\mathrm{VII}}(x)=-\frac{61+479ss+179s^4+s^6}{c^7}\Big]&&F=-\tfrac1{5040}(61+662tt+1320t^4+720t^6)
\end{array}
\]
\begin{center}etc.;\end{center}
\par\endgroup
[folglich nach Art.~19 oder 20:]
\begingroup\footnotesize
\[
\begin{aligned}
x+\Omega&=x+\tfrac12tyy-(\tfrac5{24}t+\tfrac18t^3)y^4\ldots\\
&\quad+\omega\left\{1-\tfrac12(1+tt)yy+(\tfrac5{24}+\tfrac7{12}tt+\tfrac38t^4)y^4\ldots-i\left(ty-(\tfrac56t+\tfrac12t^3)y^3\ldots\right)\right\}\\
&\quad+\omega\omega\left\{\tfrac12tyy\ldots+i\left(\tfrac12y-(\tfrac5{12}+\tfrac14tt)y^3\ldots\right)\right\}\\
&\quad+\omega^3\left\{-(\tfrac16-\tfrac13tt)yy\ldots+i(\tfrac16ty\ldots)\right\}\\
&\quad+\omega^4\left\{\ldots-i(\tfrac1{24}y\ldots)\right\}
\end{aligned}
\]
\begin{center}etc.\end{center}
\par\endgroup
\bigskip
\begin{center}
[Reihen zwischen $\varphi$, $\psi$ und $\omega$.]\par\medskip
[21.]
\end{center}
[Zur numerischen Berechnung von $\varphi$ aus $x$, Art.~1, oder von $X$ aus $\Phi$, Art.~13, dienen die folgenden Gleichungen zwischen $\varphi$ und $\psi$.

Die Länge des elliptischen Meridianbogens, dessen große Halbachse $a$ und dessen Exzentrizität $e$ ist, vom Äquator bis zur Breite $\varphi$ sei $x$. Dann gilt]
\vfill\hfill$22^*$\hspace*{2cm}\par
\endgroup
